\documentclass[notuble,10pt,a4paper]{leaflet}
\usepackage{color}
\usepackage{flowfram}
\usepackage{graphicx}
\usepackage{wrapfig}
\usepackage{hyperref}
\usepackage{rotating}
\usepackage{multirow}
\usepackage{array}
\pagestyle{empty}
\usepackage{multirow}
\usepackage{titlesec}
\usepackage[usenames,dvipsnames]{xcolor}

\newcommand*{\rom}[1]{\expandafter\@slowromancap\romannumeral #1@}
\titleformat*{\section}{\color{Blue}}
%\newenvironment{tips}
%{
%\noindent
%\begin{table}
%\begin{center}
%\begin{tabular}{m{.2cm} m{\textwidth}}
%\rotatebox{90}{\textbf{\resizebox{0.45\textwidth}{!}{FISAT}}}&
%}
%{\end{tabular}
%\end{center}
%\end{table}
%\bigskip
%}

%FONT Change
%\renewcommand{\familydefault}{cmss} 
%To Draw a horizontal Line
\newcommand{\sectionline}{
  \nointerlineskip \vspace{\baselineskip}
  \hspace{\fill}\rule{0.8\linewidth}{.7pt}\hspace{\fill}
  \par\nointerlineskip \vspace{\baselineskip}
}
%\AddToBackground{2}{\includegraphics[width=29.7cm]{bkf}}
%\AddToBackground{2}{\includegraphics[width=29.7cm]{cmy}}
% Make a border along the top of each page
\vtwotonetop{1cm}{0.6\paperwidth}{[cmyk]{0,0.91,0.92,0.41}}{topleft}%
{0.4\paperwidth}{[cmyk]{0,0.91,0.92,0.41}}{topright}

\vtwotonebottom{1cm}{0.6\paperwidth}{[cmyk]{0,0.91,0.92,0.41}}{bottomleft}%
{0.4\paperwidth}{[cmyk]{0,0.91,0.92,0.41}}{bottomright}

%\vtwotonebottom[1,6]{1cm}{0.6\paperwidth}{[cmyk]{0,0.91,0.92,0.41}}{bottomleft}%
%{0.4\paperwidth}{[cmyk]{0,0.91,0.92,0.41}}{bottomright}

\begin{document}
\begin{center}
{\large \textbf{ICEFOSS 2019}} \\[.3cm]
\textit{organized by}\\[.4cm]
{\large \textbf{FISAT Free Software Cell (FFSC)}}\\ [.3cm]
\includegraphics[scale=.3]{data/ffsc.jpg}\\
\textbf{{\large $19^{th}$ Feb - $20^{th}$ Feb 2019}}\\[.5cm]
\textit{{\small at}}\\[.3cm]
{\large\textbf{ Federal Institute of Science And Technology (FISAT)$^{\tiny{\mbox{TM}}}$}} \\[.4cm]
\includegraphics[scale=.13]{logod}\\
\vfill
Event Coordinators\\ 
\textbf{FFSC}\\
Dept. of Computer Science \& Engineering\\
Federal Institute of Science And Technology (FISAT)\\
Kerala – 683577
\end{center}
%\begin{titlepage}
%\title{
%}
%\vfill
%\date{}
%\author{
%\includegraphics[scale=0.2]{ffsc}\\
%\textit{at}\\
%Federal Institute of Science And Technology\\(FISAT)$^{\tiny{\mbox{TM}}}$ \\\\
%\includegraphics[scale=.1]{logod}\\
%}
%\end{titlepage}


%\maketitle
\thispagestyle{empty} 
\newpage
\section{{\Large{ICEFOSS-19}}}
About 7 years ago, an important initiative was taken by FISAT Free Software Cell to enhance the knowledge of FOSS and introduce modern technologies to students, faculty, colleagues and to all other FOSS enthusiasts. ICEFOSS is a conference for enthusiasts of FOSS. ICEFOSS 2019 is the VII\textsuperscript{th} national conference, which will be conducted on 19th and 20th of Feb 2019. 

\section{{\Large{Fisat Free Software Cell(FFSC)}}}
he FISAT Free Software Cell (FFSC), established in 2006, is dedicated to promote computer users` rights to use, study, copy, modify, and redistribute computer programs. The FFSC promotes the development and use of free software, particularly the GNU operating system, used widely in its GNU/Linux variant.
\section{{\Large Event Outline}}

These workshops and the women hacakathon will provide an introduction to the latest and popular programming languages, design patterns and state of the art machine learning techniques and big data analysis. Most importantly the event will be introducing Free and Open Source Softwares and technologies that can be used in a productive environment and respective domains. The event consists of:
\begin{itemize}
\item Women Hackathon
\item Go Programming
\item Introduction To Big Data And Hadoop
\item Introduction to Docker and Kubernetes
\item Introduction To Machine Learning Using Apache Spark
\end{itemize}

\newpage
\section{{\Large Women Hackathon}}
ICEFOSS 2019, scheduled on February 19 and 20, brings to you ‘Women Hackathon’, conducted in association with ICFOSS. It is a hackathon exclusively for women which aims to connect and inspire women, and be a stepping stone in furthering education and technological career. Women hackathon provides both an opportunity to create space that will make women leaders of technology.

\section{{\Large How to Apply}}
The registration for this event is done through the  offical ICEFOSS website.\newline\newline
\url{http://icefoss.fisat.ac.in/}\newline\newline
Registration is strictly online, and no other mode of applications, will be entertained. Seats will be alloted to participants who register their names and pay the registration fee of rupees 350(2-days workshop or 2-days Hackathon + lunch).
Last date for enrolment and for submission of permission letter is 18th Feb 2019, till 11am.
A list of selected participants will be put up on this website on 18th Feb 2019, afternoon. The selected participants will also be informed by email.

\emph{\textbf{LAST DATE FOR ONLINE REGISTRATION}}
\emph{\textbf{$18^{th}$ Feb 2019, till 11am}}

\section{\Large{Resource Persons}}

\begin{itemize}
\item Johny James, Senior Software Engineer at Myntra
\item Naufal Ibrahim, Free Software Consultant and Trainer.
\item Haridas N, Principal Engineer at Imaginea Technologies Inc.
\item Midhun K, Senior Software Engineer at United Layer
\item Pankaj Kumar, Assistant Professor, FISAT
\end{itemize}
 
\section{{\Large Who Should Attend}}
This workshop will benefit all students, faculty members and other FOSS enthusiasts. We encourage students who are currently pursuing computer science or information technology to participate in this event, since the speakers are industrial experts and the workshops are on cutting-edge technologies.\textbf{ We also strongly encourage women to take part in the FOSS Women Hackathon }, so as to gain experience, to collaborate with industrial experts and to apply their technical skills.

\textbf{\newline
Note A}
The number of seats are limited. Thus seats will be alloted in a "first come first serve" basis, to those who pay the registration fees before the deadline.


\section{{\Large Accommodation and other support}}

Accommodation will be provided within the college campus, on request. Students who want such accommodation are supposed to make request as soon as they finish registration, by mailing the remote center co-ordinator(jestinjoy@fisat.ac.in) or by contacting the event co-ordinators. Lunch will be provided for the participants, and no extra payment other than the registration fees will be required. For special requests, directly contact the Remote Center Coordinator - Mr. Jestin Joy (9446764386).

\newpage
\section{{\Large Eligibility}}

\begin{itemize}
\item Teachers employed in engineering colleges are eligible to attend this workshop. This includes teachers who are visiting faculty, and contract employees.
\item All Students (current/pass-outs) can participate.
\item FOSS enthusiasts and others can also participate.
\end{itemize}

\section{{\Large Duration and Venue}}
The ICEFOSS 2019 workshops on Go Programming, Introduction To Big Data And Hadoop, Introduction to Docker and Kubernetes, ML using Apache Spark and the two day FOSS Women Hackathon will be conducted simultaneously in multiple labs at the college.The workshop will include the delivery of interactive lectures from eminent speakers, locally organized tutorials and lab sessions.\\
Each of the workshops will be conducted full day, from sharp 9:15am to 4:15pm.The women hackathon will be a two day event(accommodation can be provided), and requires individual participants, rather than a team.

\section{{\Large Event Registration Fee}}

Only a minimal amount of \textbf{rupees 350}  is required for attending the two days workshop/women hacakathon plus the lunch. Registration should be done via the official ICEFOSS website.
\newpage

%\begin{center}
%\centering
%\includegraphics[scale=0.4]{ffsc}
%\end{center}
%\begin{center}
%\begin{tips}
%\begin{tabular}{l}
%{\resizebox{5cm}{!}{ONE}}\\ 
%{\resizebox{5cm}{!}{MILLION}}\\ 
%{\resizebox{5cm}{!}{CODES TO}} \\ 
%{\resizebox{5cm}{!}{THE COMMUNITY}}\
%\end{tabular} 
%\end{tips}
%\end{center}
%\begin{center}



\section{{\Large Event Center}}
\textbf{Federal Institute of Science And Technology (FISAT)}, Angamaly, Kerala has been approved as one of the remote centre of this project. FISAT was established in 2002 under the aegis of Federal Bank Officers' Association Educational Society (FBOAES). The college has carved a niche for itself in education world, eloquently demonstrated by the flying colors attained by its students in academics, placements as well as in extra curricular and co-curricular activities. The campus is a class apart with magnificent infrastructure and has an ambiance that kindles creativity and encourages innovation.

The college has established research centres like Centre for High Performance Computing (CHPC), Instrumentation Research and Consultancy Centre (IRACC) and Centre for Research and Innovation in Signal Processing (CRISP) to nurture research and development. Under research promotion scheme, AICTE funded research centres were sanctioned for Signal Processing in ECE department and High Performance Computing in CSE department. 

For more details, visit \url{www.fisat.ac.in}

\section{{\Large Address for Communication}}
\textbf{Yedhin Kizhakkethara}\\
FFSC Core, Department of CSE,\\ 
FISAT\\
Tel.: +91-9497656739\\
Email:  \url{yedhin1998@gmail.com}

\section{{\Large Event Center Co-ordinator}}
\textbf{Mr. Jestin Joy}\\
Federal Institute of Science And Technology (FISAT)\\
Hormis Nagar, Mokkanoor, Angamaly\\
Mob: +91 9446764386\\
E-mail:\url{jestinjoy@fisat.ac.in}\\

\vfill
\footnotesize{\textcopyright FISAT | Prepared using \LaTeXe}



%\end{center}
\end{document}
